\section*{Hoofdstuk \ref{ch:b3:radicals-carbenes} --- Radicalen, carbenen en omleggingen}

\begin{solution}{exo:b3:radicals-carbenes:1}
Propaan heeft zes primaire en twee secundaire waterstofatomen. Chlorering: $6 \times 1 : 2 \times 3.9 = 6 : 7.8$,
d.w.z.\ 43\,\% 1-chloorpropaan en 57\,\% 2-chloorpropaan. Bromering: $6 : 164$, 4\,\% en
96\,\%.
\end{solution}

\begin{solution}{exo:b3:radicals-carbenes:2}
$\ce{CH2}$: triplet; $\ce{CCl2}$: singulet (de vrije elektronenparen van chloor verhogen de lege orbitaal p).
Dichloorcarbeen addeert stereospecifiek: \textit{trans}-1,1-dichloor-2,3-dimethylcyclopropaan.
\end{solution}

\begin{solution}{exo:b3:radicals-carbenes:3}
\textit{cis}-1,2-Diethylcyclopropaan: de \omterm{def:b3:radicals-carbenes:carbenoid}{carbenoïde} addeert in één stap aan één vlak.
\end{solution}

\begin{solution}{exo:b3:radicals-carbenes:4}
Cyclohexanon: $\varepsilon$-caprolacton (oxepan-2-on). Acetofenon: fenylacetaat,
$\ce{CH3CO2C6H5}$; fenyl migreert bij voorkeur boven methyl.
\end{solution}

\begin{solution}{exo:b3:radicals-carbenes:5}
Negen primaire en één tertiair waterstofatoom. Chlorering: $9 : 5.1$, 64\,\% 1-chloor-2-methylpropaan
en 36\,\% 2-chloor-2-methylpropaan. Bromering: $9 : 1600$, 0.6\,\% en 99.4\,\%.
\end{solution}

\begin{solution}{exo:b3:radicals-carbenes:6}
$k_c = 3.0 \times k_H[\mathrm{Bu_3SnH}] = 3.0 \times \num{2.4e6} \times 0.020 = \qty{1.4e5}{s^{-1}}$.
\end{solution}

\begin{solution}{exo:b3:radicals-carbenes:7}
Aanval op C5: 5-exo-trig, begunstigd; op C6: 6-endo-trig, volgens de
regels begunstigd maar hier trager. 4-exo-tet: begunstigd. 5-endo-trig:
benadeeld. 5-exo-dig: begunstigd.
3-exo-dig: benadeeld. 5-endo-dig: begunstigd.
\end{solution}

\begin{solution}{exo:b3:radicals-carbenes:8}
Verlies van water geeft een tertiair, benzylisch kation; op het naburige
koolstofatoom kunnen een fenyl- en een methylgroep verhuizen, en fenyl,
met de hogere \omterm{def:b3:radicals-carbenes:migratory}{migratiebereidheid}, migreert. Het door zuurstof
gestabiliseerde kation verliest een proton: 3,3-difenylbutan-2-on.
\end{solution}

\begin{solution}{exo:b3:radicals-carbenes:9}
Curtius: $\ce{C6H5CON3 -> C6H5NCO + N2}$. Met water verliest het carbaminezuur $\ce{CO2}$: aniline
(er ontstaat wat difenylureum wanneer aniline meer \omterm{def:b3:radicals-carbenes:curtius}{isocyanaat} ontmoet). Met
ethanol: ethyl-\textit{N}-fenylcarbamaat, $\ce{C6H5NHCO2C2H5}$.
\end{solution}

\begin{solution}{exo:b3:radicals-carbenes:10}
Chloor: primair $421.7 - 432 = \qty{-10}{kJ/mol}$, tertiair $405.4 - 432 = \qty{-27}{kJ/mol}$; broom: $+56$ en
$+39$. Beide onttrekkingen door chloor zijn exotherm, met vroege
overgangstoestanden die weinig voelen van het verschil van
\qty{16}{kJ/mol}; beide onttrekkingen door broom zijn endotherm, met late
overgangstoestanden waarvan de activeringsenergieën het grootste deel
ervan verschillen:
$\eu^{16000/(8.314 \times 298)} \approx 600$, de orde van de waargenomen verhouding.
\end{solution}

\begin{solution}{exo:b3:radicals-carbenes:11}
Oxim \textit{E} (OH anti ten opzichte van C2, dat de methylgroep draagt): C2 migreert, stikstof komt
ernaast: 7-methylazepan-2-on. Oxim \textit{Z}: C6 migreert: 3-methylazepan-2-on. Bij de
\omterm{def:b3:radicals-carbenes:beckmann}{beckmannomlegging} beslist de geometrie van het oxim; bij de
\omterm{def:b3:radicals-carbenes:baeyer}{baeyer-villigeroxidatie} de \omterm{def:b3:radicals-carbenes:migratory}{migratiebereidheid}, en het secundaire
koolstofatoom verhuist:
7-methyloxepan-2-on.
\end{solution}

\begin{solution}{exo:b3:radicals-carbenes:12}
Gecycliseerde fractie $k_c/(k_c + k_H c)$: 0.91, 0.49 en 0.09. Voor 95\,\%: $k_Hc = k_c(1/0.95 - 1)$, dus
$c = 0.0526 \times \num{2.3e5}/\num{2.4e6} = \qty{5.0e-3}{mol/L}$. Men houdt die aan door het hydride langzaam toe te
voegen met een spuitpomp, of door een katalytische hoeveelheid
tinhalogenide te gebruiken met een reductor die het hydride regenereert.
\end{solution}

\begin{solution}{pb:b3:radicals-carbenes:1}
\textbf{1.} $\ce{C6H10O + NH2OH -> C6H11NO + H2O}$.

\textbf{2.} Er is zuur nodig om de hydroxygroep van het additie-intermediair
als water te verwijderen, maar te veel zuur protoneert het
hydroxylamine, dat dan niet meer addeert.

\textbf{3.} Nee: de twee koolstofatomen die aan het koolstofatoom van C=N gebonden zijn, zijn door de
symmetrie van de ring gelijkwaardig.

\textbf{4.} In het oxim \textit{E} wijst de OH weg van het koolstofatoom C2 met de methylgroep (anti ten opzichte
ervan), in het oxim \textit{Z} ernaartoe.

\textbf{5.} Inversie aan een oximstikstofatoom dat een elektronegatief
zuurstofatoom draagt, heeft een hoge barrière.

\textbf{6.} Het protoneert (of sulfoneert) de hydroxygroep, waardoor die
een goede uittredende groep wordt, water.

\textbf{7.} De binding C--C van de ring die anti staat ten opzichte van de
uittredende groep, verhuist van koolstof naar stikstof, dat daardoor in
de ring ingevoegd wordt: zes atomen worden er zeven.

\textbf{8.} Een cyclisch nitriliumion.

\textbf{9.} $\ce{C6H11NO -> C6H11NO}$: een isomerisatie, van oxim tot lactam.

\textbf{10.} Caprolactam, azepan-2-on, het zevenledige cyclische amide van
6-aminohexaanzuur; zijn ringopeningspolymerisatie geeft nylon-6.

\textbf{11.} \qty{98.0}{g/mol} en \qty{113.0}{g/mol}.

\textbf{12.} $\qty{1.0e6}{g}/\qty{98.0}{g/mol} = \qty{10204}{mol}$; $\times 0.98 \times 0.98 = \qty{9800}{mol}$; $\times \qty{113.0}{g/mol} = \qty{1.11e6}{g}$, \qty{1.11}{t}.

\textbf{13.} Oxim: $\num{10204} \times 0.98 = \qty{1.00e4}{mol}$; $\ce{H2SO4}$: $\qty{1.30e4}{mol}$, wat evenveel mol
$\ce{(NH4)2SO4}$ geeft: $\qty{1.30e4}{mol} \times \qty{132.1}{g/mol} = \qty{1.7}{t}$, meer dan het caprolactam.

\textbf{14.} $\ce{C6H10O + RCO3H -> C6H10O2 + RCOOH}$.

\textbf{15.} Het tetraëdrische adduct van het peroxyzuur op het
carbonylkoolstofatoom (het intermediair van Criegee), een
peroxyhemiketal.

\textbf{16.} Een groep $\ce{CH2}$ van de ring (de twee zijn gelijkwaardig) verhuist naar het dichtstbijzijnde
peroxidezuurstofatoom; het carboxylaat $\ce{RCOO-}$ vertrekt.

\textbf{17.} $\varepsilon$-Caprolacton, oxepan-2-on; zijn ringopeningspolymerisatie geeft
poly($\varepsilon$-caprolacton), een biologisch afbreekbare polyester.

\textbf{18.} Het enige bijproduct is water, terwijl een peroxyzuur een mol
carbonzuur per mol lacton achterlaat en zelf gevaarlijk is om op te
slaan.

\textbf{19.} Het secundaire koolstofatoom C2 migreert: 7-methyloxepan-2-on.

\textbf{20.} Ja: het migrerende koolstofatoom behoudt zijn configuratie.

\textbf{21.} 7-Methylazepan-2-on (C2 staat anti en migreert).

\textbf{22.} 3-Methylazepan-2-on (C6 migreert).

\textbf{23.} Beckmann: de geometrie van het oxim (de anti-groep verhuist).
Baeyer-Villiger: de \omterm{def:b3:radicals-carbenes:migratory}{migratiebereidheid} (het sterker gesubstitueerde
koolstofatoom verhuist).

\textbf{24.} Zijn twee koolstofatomen $\alpha$ zijn gelijkwaardig: één oxim, één lactam, één lacton.

\textbf{25.} Eén ton cyclohexanon geeft ongeveer \qty{1.11}{t} caprolactam.
\end{solution}
